\section{아티팩트 산출물과 완료 판정}
\label{sec:artifact-plan}

\subsection{작업 순서}

이 장은 \cref{sec:methodology}의 증명 절차를 실제 산출물과 완료 판정으로
구체화한다.
아래 순서는 7일 마감에 맞춘 \textbf{GO-MINCORE-7D} 실행 순서이면서 theorem
dependency 순서이다. Week~15의 fixed all-6 정리는 종료한 실행의 실패를
배제하는 Hoare 부분정확성이지 termination/losslessness나 non-vacuous 실행
증명이 아니다. 이 미완료 항목과 \(h\) byte codec은 이번 일정에서 보류한다.
대신 KeyGen이 제공하는 키 관계, Sign의 accepted core 계산, Verify의 bounded
core predicate를 증명하고 decoded signature object 경계에서 합성한다. 앞 단계가
닫히지 않으면 후속 headline theorem에 결론을 새 가정으로 숨기지 않고 정확한
focused actual 경계에서 미완료로 남긴다.

\begin{longtable}{L{13mm}L{37mm}L{92mm}}
\toprule
단계 & 산출물 & 닫아야 할 내용 \\
\midrule
\endhead
1일차 & 범위와 baseline 동결 &
actual procedure, mode-2 상수, theorem signature와 현재 aggregate fresh
baseline을 고정한다. \\
2일차 & KeyGen headline theorem &
actual matrix/finalization 경계에서 \textup{(KG-1)}--\textup{(KG-4)}와
\(As=qj\pmod{2q}\)를 도출한다. accepted parent lift는 기존 caller 사실로
즉시 합성되는 경우에만 포함한다. \\
3--4일차 & Sign headline theorem &
\textup{(S-1)}--\textup{(S-7)}과 실제 두 norm predicate를 세 core
procedure의 순차 호출로 합성한다. \\
5--6일차 & Verify 및 제한 합성 &
\textup{(V-1)}--\textup{(V-8)}과 \(reject=0\) 동치를 bounded actual core
경계에서 증명하고, 동일 key/message transcript와 decoded-object identity 아래
accepted Sign 결과가 Verify core에서 승인됨을 증명한다. \\
7일차 & Paper artifact &
source hash, extraction identity, safety evidence, theorem manifest, fresh
\proc{-no-eco} build, proof-hole/axiom scan, 독립 audit와 TCB를 제공한다. \\
\bottomrule
\end{longtable}

\subsection{각 headline theorem의 완료 기준}

\begin{acceptancecriterion}[기계검증 완료]
KeyGen, Sign, Verify 정리는 각각 다음 조건을 모두 만족할 때만 완료로 판정한다.
\begin{enumerate}
  \item 최종 정리가 pinned extraction의 실제 프로시저를 직접 호출한다.
  \item 성공, decoded equality, trace equality 또는 최종 predicate를 전제로
        넣지 않고 실제 하위 postcondition에서 도출한다.
  \item 모든 representation, bounds, frame 및 memory-disjointness 전제를
        theorem statement에 공개한다.
  \item Hoare partial correctness, losslessness, termination을 구분하여 표기한다.
  \item 신규 타깃 개별 fresh compile과 aggregate fresh \proc{-no-eco}
        compile이 모두 통과한다.
  \item proof hole, authored axiom, debug declaration, extraction/table drift가 없다.
  \item claim ledger, theorem graph, proof manifest와 논문의 상태가 일치한다.
\end{enumerate}
\end{acceptancecriterion}

\subsection{안전성 및 TCB}

Jasmin compiler와 EasyCrypt extraction의 의미 연결은 safe program에 대해
적용된다 \cite{jasmin-extraction}. 따라서 assembly-level correctness를
주장하려면 논문 대상 프로시저의 memory safety 증거를 함께 제공한다. 이를
제공하지 못한 프로시저는 ``extracted EasyCrypt model 수준''의 결과로만
표현한다.

TCB에는 적어도 다음을 명시한다.

\begin{itemize}
  \item pinned HAETAE 명세와 그 EasyCrypt formalization,
  \item Jasmin instruction semantics, compiler proof 및 \proc{jasmin2ec},
  \item EasyCrypt kernel/SMT backends와 사용한 버전,
  \item generated mode-2 tables와 extraction scripts,
  \item theorem statement와 구현 배열을 수학 object로 해석하는 bridge.
\end{itemize}

\subsection{명시적 비주장}

최종 핵심 연산 논문은 다음을 주장하지 않는다.

\begin{itemize}
  \item KeyGen sampler와 hyperball sampler의 정확한 출력분포,
  \item KeyGen/Sign rejection loop의 일반 termination 또는 기대 반복 횟수,
  \item 모든 canonical HBZ stream의 rANS encoder 성공,
  \item \(h\) codec, metadata, padding을 포함한 full 1474-byte canonical
        parser,
  \item 임의 malformed signature의 완전한 rejection characterization,
  \item 전체 public API의 분포 동등성,
  \item HAETAE 구현의 EUF-CMA 또는 논문 보안 환원 전체.
\end{itemize}

\subsection{예상 작업량과 보고 제목}

현재 추출, memory bridge, raw \(\mu\) transcript, signature prefix 및
HBZ/rANS vertical slice를 재사용한다 \cite{claim-ledger,week14-report}.
활성 산출물의 시간 한도는 7일이다. 이는 주장 강도를 높이는 근거가 아니라
범위를 decoded-object MINCORE로 동결하는 제약이다. upper caller lift가 일정에
맞지 않으면 관찰용 모델로 대체하지 않고 focused actual core 경계를 명시한다.
full byte codec, public API, termination과 분포까지 포함하는 기존 장기 계획은
이번 산출물 범위 밖이다.

최종 보고 제목은 다음과 같이 고정한다.

\begin{center}
\bfseries
HAETAE-2 Jasmin 구현의 KeyGen 키 방정식,\\
Sign 응답·힌트 계산 및 Verify 재구성·승인 계산에 대한 기계검증
\end{center}

``HAETAE 전체 기능정확성 검증'' 또는 ``HAETAE 구현 보안 증명''이라는 제목은
위 비주장들이 닫히기 전에는 사용하지 않는다.
